\justifying
\NSFCBodyText

\subsubsection{年度研究计划}

以下按四个执行年度编排，具体起止年月须以当年度项目指南和批准通知为准。

\paragraph{第一执行年度：操作定义、数据治理与操纵验证。}完成两个候选公开中文来源的许可、版本、隐私和独立 token 审计，确定来源 A/B；冻结一类确认性中文语体和一类稳健性复核现象，形成结构完整性、语义保持、事实一致与语体属性的盲法标注规范。实现困惑度型和来源/规范性型两类质量表征，构建分数重叠样本层与六臂反事实语料生成器。完成每臂一次、单次暂定 0.10B token 的非推断性工程先导，检验评分重叠、人工一致性、sham 有效性、语义保持、token/packing 平衡和训练稳定性；单次先导不用于估计种子间方差或证明等效。里程碑为数据卡、标注一致性报告、六臂平衡表、冻结的统计分析计划和可复现 smoke test；任一操纵不合格先返工，不进入确认性训练。

\paragraph{第二执行年度：来源 A 的双路径确认与动力学检验。}在固定约 1B 开源基座上执行六个处理臂、六个全新成对主种子的继续预训练，单次 token 暂定 0.40B，并以首年吞吐、独立语料规模和功效模拟校准。完成通用语言/知识与目标语体/长尾两个主能力终点的运行级分析，检验结构完整相对破损与 sham 的恢复特异性，以及同语体非冗余补回相对主流和窄支持补回的支持特异性。同步计算能力定向梯度对齐，判断其是否满足处理先行、恢复共同逆转和留出预测增益。里程碑为两条路径的效应/等效区间、候选机制判断和继续/收缩/停止决定；失败分支不进入后续外推。

\paragraph{第三执行年度：独立来源的冻结协议复核。}将来源 A 中存活路径的处理规则、主终点、主动力学读出、效应方向、预测区间和代码冻结，在来源 B 上复核结构完整/破损和主流/同语体非冗余补回四个关键臂，暂定每臂六个成对种子、每次 0.40B token。该阶段只运输冻结主对比，不宣称再次验证 sham 或窄支持特异性；先区分来源准入失败、操纵失败、估计不精确与真实效应异质，再评价方向、冻结区间覆盖以及一次固定预算校准的误差变化。里程碑为有限运输复核报告和“直接运输、条件校准、证据不足、拒绝外推”判定，不新增来源追逐一致性。

\paragraph{第四执行年度：状态/尺度边界探索与综合凝练。}仅对前三年通过闸门的路径开展处理前知识状态随机分叉；若至少一条路径在来源 B 保持方向，再以同 tokenizer 家族的较大模型完成少数关键臂探索。两类实验均只报告方向、区间和失败模式，不作普遍显著性或跨架构因果主张。整合路径效应、能力定向梯度读出、来源边界与替代解释，完成论文、可复现代码、数据处理协议和可公开派生标注的整理。里程碑为双路径机制的支持/部分支持/反证结论及其适用范围。

\subsubsection{预期研究结果}

（1）形成一套可复核的中文语料结构完整性、模型熟悉度与静态质量表征分层方法，以及由结构修复、sham、主流补回、同语体非冗余补回和同语体窄支持补回构成的反事实协议。该产物强调操作定义和审计证据，不追求成为通用质量评分标准。

（2）获得结构缺陷削减相对破损与 sham 的运行级效应和反证结果，评估能力定向梯度对齐是否满足处理先行、恢复共同逆转与留出预测增益，并给出该结构效应在独立中文来源中的运输、校准或拒绝外推边界。

（3）获得非冗余语义支持损失、补回恢复和支持广度的运行级效应及等效区间，区分语体曝光与非冗余广度，检验能力定向梯度读出，并形成状态—来源—尺度边界图。具体论文数量、投稿期刊/会议、数据开放范围和软件版本根据真实结果与语料许可确定，不把平台、标准、商业模型或全部原始语料公开列为必达成果。
